\documentclass[11pt]{article}

\usepackage[margin=1in]{geometry}
\usepackage{amsmath,amssymb,amsthm}
\usepackage{hyperref}
\usepackage{microtype}

\title{\textbf{BSD Curve-wise Spectral Leap}}
\author{Inacio F.\ Vasquez \\ Independent Researcher}
\date{}

% --- Theorem Environments ---
\theoremstyle{plain}
\newtheorem{theorem}{Theorem}
\newtheorem{lemma}[theorem]{Lemma}

\theoremstyle{remark}
\newtheorem{remark}[theorem]{Remark}

% --- Notation ---
\newcommand{\Q}{\mathbb{Q}}
\newcommand{\Z}{\mathbb{Z}}
\newcommand{\hath}{\hat h}

\begin{document}

\maketitle

\begin{center}
\textbf{Status summary.}
\begin{itemize}
\item Curve-wise spectral leap: \textbf{Proven for each curve}, uniformity not assumed.
\item Dependencies: Mordell--Weil theorem, positive definiteness of Néron--Tate height, lattice minimum.
\item Lean references: \texttt{URF\_Core.lean → BSD.*}
\end{itemize}
\end{center}

\tableofcontents

\section{Néron--Tate Form and Positive Definiteness}

\begin{definition}[Néron--Tate height quadratic form]
Let $E/\Q$ be an elliptic curve. The Néron--Tate quadratic form is
\[
\hath(P) : E(\Q) \to \R.
\]
\end{definition}

\begin{lemma}[Positive definiteness]\label{lem:bsd-positive}
The quadratic form $\hath$ is positive definite on $E(\Q)/E(\Q)_{\mathrm{tors}}$.
\end{lemma}

\begin{proof}
$\hath$ is quadratic on $E(\Q)\otimes_\Z \R$.  
Kernel is torsion points.  
Hence it descends to a positive definite form on the non-torsion lattice.  
\textbf{Lean stub reference:} \texttt{URF\_Core.lean → BSD.positive\_definite}
\end{proof}

\section{Lattice Minimum}

\begin{lemma}[Lattice minimum]\label{lem:bsd-lattice}
For a positive definite form on the lattice $E(\Q)/E(\Q)_{\mathrm{tors}}$, the minimum over nonzero points is strictly positive.
\end{lemma}

\begin{proof}
Any positive definite quadratic form on a lattice $\Z^r$ attains a strictly positive minimum on nonzero lattice points.  
\textbf{Lean stub reference:} \texttt{URF\_Core.lean → BSD.lattice\_minimum}
\end{proof}

\section{Curve-wise Spectral Gap}

\begin{theorem}[Curve-wise spectral gap]\label{thm:bsd-curvewise}
For each fixed $E/\Q$, there exists $\lambda(E)>0$ such that
\[
\hath(P) \ge \lambda(E) \quad \forall P \in E(\Q)\setminus E(\Q)_{\mathrm{tors}}.
\]
\end{theorem}

\begin{proof}
Apply Lemmas \ref{lem:bsd-positive} and \ref{lem:bsd-lattice}.  
This gives a strictly positive minimum for non-torsion points, yielding a spectral gap.  
\textbf{Lean stub reference:} \texttt{URF\_Core.lean → BSD.curvewise\_gap}
\end{proof}

\begin{remark}[Scope]
The result is curve-wise; no uniform lower bound across all curves is assumed.  
\textbf{Frozen module:} Yes.
\end{remark}

\end{document}

